\documentclass[12pt]{ctexart}
\usepackage[a4paper,margin=2.2cm]{geometry}
\usepackage{amsmath,amssymb,amsthm}
\usepackage{bm}
\usepackage{enumitem}
\usepackage{hyperref}
\usepackage{fancyhdr}
\usepackage{tikz}
\usetikzlibrary{arrows.meta,positioning,calc}
\setlength{\headheight}{15pt}
\setCJKfamilyfont{handfont}{Kaiti SC}

\hypersetup{
  colorlinks=true,
  linkcolor=black,
  urlcolor=blue
}

\pagestyle{fancy}
\fancyhf{}
\lhead{高数基础}
\chead{第三章 中值定理及导数应用}
\rhead{\thepage}

\renewcommand{\thesection}{\S\arabic{section}.}
\renewcommand{\thesubsection}{\arabic{subsection}.}
\renewcommand{\thesubsubsection}{(\arabic{subsubsection})}

\newtheorem{theorem}{定理}[section]
\newtheorem{example}{例}[section]
\theoremstyle{definition}
\newtheorem{definition}{定义}[section]
\newtheorem{remark}{注}[section]

\newcommand{\hand}[1]{{\CJKfamily{handfont}#1}}

\title{\textbf{高数基础\\第三章\quad 中值定理及导数应用}}
\author{根据 PDF 讲义整理}
\date{\today}

\begin{document}
\maketitle
\tableofcontents
\newpage

\section{微分中值定理}

\subsection{前言}
\paragraph{(1) 解决问题}
中值定理主要解决开区间内\textbf{有导数的存在性}问题。

\paragraph{(2) 知识体系}
\begin{figure}[htbp]
\centering
\begin{tikzpicture}[
  node distance=16mm and 18mm,
  box/.style={draw, rounded corners=8pt, align=center, inner sep=7pt, font=\small},
  arrow/.style={-{Latex[length=2.4mm]}, thick}
]
\node[box] (fermat) {费马引理};
\node[box, below=of fermat] (rolle) {罗尔中值定理};
\node[box, right=32mm of rolle] (lagrange) {拉格朗日中值定理};
\node[box, below=of lagrange] (cauchy) {柯西中值定理};

\node[right=6mm of fermat, align=left, font=\small] (fermat-note) {去掉定点时界\\证明极值点};
\node[right=6mm of rolle, align=left, font=\small] (rolle-note) {去掉定值时界\\证明驻点};
\node[right=6mm of lagrange, align=left, font=\small] (lagrange-note) {去掉定值时界\\证明导数};
\node[right=6mm of cauchy, align=left, font=\small] (cauchy-note) {去掉定值时界\\求值};

\draw[arrow] (fermat) -- (rolle);
\draw[arrow] (rolle) -- (lagrange);
\draw[arrow] (lagrange) -- (cauchy);
\end{tikzpicture}
\caption{知识体系图}
\end{figure}

\subsection{极值的定义}
\begin{definition}
若存在 $\delta>0$，使得对任意
\[
x\in (x_0-\delta,x_0+\delta),\quad x\neq x_0
\]
恒有
\[
f(x)>f(x_0)\quad [f(x)<f(x_0)],
\]
则称 $x=x_0$ 是函数 $y=f(x)$ 的极小（大）值点；$f(x_0)$ 是函数的极小（大）值。
\end{definition}

\begin{remark}
\begin{enumerate}[label=\arabic*)]
  \item 极值点一定在区间内部取得；
  \item 极大值点、极小值点统称为极值点；极大值、极小值统称为极值。
\end{enumerate}
\end{remark}

\begin{figure}[htbp]
\centering
\begin{tikzpicture}[scale=0.9, >=Latex]
  \draw[->] (0,0) -- (5.2,0) node[right] {$x$};
  \draw[->] (0,0) -- (0,3.8) node[above] {$y$};
  \draw[thick, smooth] plot coordinates {(0.6,1.2) (1.5,2.5) (3,1.1) (4.5,2.8)};
  \fill (1.5,2.5) circle (2pt) node[above] {$B$};
  \fill (3,1.1) circle (2pt) node[below] {$C$};
  \fill (4.5,2.8) circle (2pt) node[above] {$D$};
  \node[below] at (0.6,0) {$A$};
\end{tikzpicture}
\caption{$B,C$ 为极值点}
\end{figure}

\subsection{极值点处的可导性情况}
\paragraph{（1）定理 1}
若 $x=x_0$ 是函数 $y=f(x)$ 的极值点，则
\[
f'(x_0)=0 \quad \text{或} \quad f'(x_0)\ \text{不存在}.
\]

\paragraph{（2）定理 2：费马定理}
若 $y=f(x)$ 在点 $x=x_0$ 处取得极值且可导，则
\[
f'(x_0)=0.
\]

\begin{remark}
本定理可理解为“可导函数在极值点的必要条件”。
\end{remark}

\begin{figure}[htbp]
\centering
\begin{tikzpicture}[scale=0.9, >=Latex]
  \draw[->] (0,0) -- (5.1,0) node[right] {$x$};
  \draw[->] (0,0) -- (0,3.5) node[above] {$y$};
  \draw[dashed] (0.7,2.6) -- (4.5,2.6);
  \draw[dashed] (1.9,1.25) -- (3.7,1.25);
  \draw[thick, smooth] plot coordinates {(0.8,2.2) (1.4,2.8) (2.4,1.9) (3.3,0.9) (4.4,2.4)};
\end{tikzpicture}
\caption{极值点处的导数必要条件示意图}
\end{figure}

\hand{即上述定理反之不成立。}

\begin{example}[反例 1]
设
\[
y=x^3.
\]
此时
\[
y'=3x^2,\qquad y'(0)=0,
\]
\hand{但 $x=0$ 不是函数 $y=x^3$ 的极值点。}
\end{example}

\begin{figure}[htbp]
\centering
\begin{tikzpicture}[scale=1.0, >=Latex]
  \draw[->] (-2.2,0) -- (2.2,0) node[right] {$x$};
  \draw[->] (0,-2.2) -- (0,2.2) node[above] {$y$};
  \draw[thick, domain=-1.2:1.2, smooth, samples=120] plot (\x,{1.1*(\x)^3});
\end{tikzpicture}
\caption{$y=x^3$ 在原点处导数为 $0$，但无极值}
\end{figure}

\begin{example}[反例 2]
设
\[
y=
\begin{cases}
2x, & x<0,\\
x, & x\ge 0.
\end{cases}
\]
\hand{函数在 $x=0$ 处不可导，但 $x=0$ 也不是函数的极值点。}
\end{example}

\begin{figure}[htbp]
\centering
\begin{tikzpicture}[scale=0.85, >=Latex]
  \draw[->] (-2.5,0) -- (2.5,0) node[right] {$x$};
  \draw[->] (0,-3.2) -- (0,2.5) node[above] {$y$};
  \draw[thick] (-1.4,-2.8) -- (0,0);
  \draw[thick] (0,0) -- (2,2);
\end{tikzpicture}
\caption{不可导并不推出极值}
\end{figure}

\subsection{罗尔中值定理}
设函数 $f(x)$ 满足
\begin{enumerate}[label=\arabic*)]
  \item 在闭区间 $[a,b]$ 上连续；
  \item 在开区间 $(a,b)$ 内可导；
  \item $f(a)=f(b)$。
\end{enumerate}
则至少存在一点 $\xi\in(a,b)$，使得
\[
f'(\xi)=0.
\]

\hand{证明可分两种情形：若函数是常数函数，则任取内点均有导数为 $0$；若不是常数函数，则由最大最小值定理与端点等值条件，可知开区间内至少有一个极值点，再由费马定理得证。}

\begin{figure}[htbp]
\centering
\begin{tikzpicture}[scale=0.95, >=Latex]
  \draw[->] (0,0) -- (5.5,0) node[right] {$x$};
  \draw[->] (0,0) -- (0,3.8) node[above] {$y$};
  \draw[dashed, red] (0.8,0) -- (0.8,2.7);
  \draw[dashed, red] (4.6,0) -- (4.6,2.7);
  \draw[dashed, red] (0.8,2.7) -- (4.6,2.7);
  \draw[dashed, red] (3.2,0) -- (3.2,1.25);
  \draw[dashed, red] (0.8,1.25) -- (4.6,1.25);
  \draw[thick, smooth] plot coordinates {(0.8,2.7) (1.6,3.05) (2.5,2.1) (3.2,1.25) (4.6,2.7)};
  \node[below, red] at (0.8,0) {$a$};
  \node[below, red] at (3.2,0) {$\xi$};
  \node[below, red] at (4.6,0) {$b$};
\end{tikzpicture}
\caption{罗尔中值定理示意图}
\end{figure}

\begin{example}[3.1]
设函数 $f(x)$ 在区间 $[a,b]$ 上连续，其中 $0<a<b$；在开区间 $(a,b)$ 内可导，且
\[
af(b)=bf(a),
\]
证明：存在 $\xi\in(a,b)$，使得
\[
f'(\xi)=\frac{f(\xi)}{\xi}.
\]

\hand{构造辅助函数}
\[
F(x)=\frac{f(x)}{x}.
\]
\hand{由题设有 $F(a)=F(b)$，再由罗尔定理知存在 $\xi\in(a,b)$，使}
\[
F'(\xi)=\frac{\xi f'(\xi)-f(\xi)}{\xi^2}=0.
\]
\end{example}

\begin{example}[3.2]
设函数 $f(x)$ 在闭区间 $[0,\pi]$ 上连续，在开区间 $(0,\pi)$ 内可导，证明：存在
\[
\xi\in(0,\pi),
\]
使得
\[
f'(\xi)=-f(\xi)\cot\xi.
\]

\hand{构造}
\[
F(x)=f(x)\sin x.
\]
\hand{若再满足 $F(0)=F(\pi)=0$，由罗尔定理得某个 $\xi\in(0,\pi)$ 满足}
\[
F'(\xi)=f'(\xi)\sin\xi+f(\xi)\cos\xi=0.
\]
\end{example}

\subsection{拉格朗日中值定理}
设函数 $f(x)$ 满足
\begin{enumerate}[label=\arabic*)]
  \item 在闭区间 $[a,b]$ 上连续；
  \item 在开区间 $(a,b)$ 内可导。
\end{enumerate}
则在 $(a,b)$ 内至少存在一点 $\xi\in(a,b)$，使得
\[
\frac{f(b)-f(a)}{b-a}=f'(\xi).
\]

\hand{该式还可写成}
\[
f(b)-f(a)=f'(\xi)(b-a).
\]
\hand{若令}
\[
\theta=\frac{\xi-a}{b-a},\qquad 0<\theta<1,
\]
\hand{则}
\[
\xi=a+\theta(b-a),\qquad f(b)-f(a)=f'(a+\theta(b-a))(b-a).
\]

\begin{figure}[htbp]
\centering
\begin{tikzpicture}[scale=0.95, >=Latex]
  \draw[->] (0,0) -- (6.2,0) node[right] {$x$};
  \draw[->] (0,0) -- (0,4.3) node[above] {$y$};
  \draw[thick, smooth] plot coordinates {(0.8,1.5) (1.6,1.9) (2.7,2.2) (3.5,3.4) (5.2,3.7)};
  \draw[thick] (0.8,1.5) -- (5.2,3.7);
  \draw[thick] (2.1,1.95) -- (4.3,3.05);
  \draw[dashed] (0.8,0) -- (0.8,1.5);
  \draw[dashed] (2.8,0) -- (2.8,2.55);
  \draw[dashed] (5.2,0) -- (5.2,3.7);
  \node[below] at (0.8,0) {$a$};
  \node[below] at (2.8,0) {$\xi$};
  \node[below] at (5.2,0) {$b$};
\end{tikzpicture}
\caption{拉格朗日中值定理示意图}
\end{figure}

\begin{remark}
\begin{enumerate}[label=\arabic*)]
  \item 可看作“割线斜率等于某点切线斜率”；
  \item 作用：证明恒等式、估值、不等式、求极限，以及中值和极值性质的推广。
\end{enumerate}
\end{remark}

\begin{example}[3.3]
设 $f(x)$ 在 $[a,b]$ 上连续，在 $(a,b)$ 内可导，证明在 $(a,b)$ 内至少存在一点 $\xi$，使
\[
\frac{bf(b)-af(a)}{b-a}=f(\xi)+\xi f'(\xi).
\]

\hand{构造}
\[
F(x)=xf(x),
\]
\hand{因为}
\[
F'(x)=f(x)+xf'(x),\qquad F(b)-F(a)=bf(b)-af(a),
\]
\hand{所以由拉格朗日中值定理，存在 $\xi\in(a,b)$ 使}
\[
\frac{F(b)-F(a)}{b-a}=F'(\xi)=f(\xi)+\xi f'(\xi).
\]
\end{example}

\begin{example}[3.4]
证明当 $x>0$ 时，
\[
\frac{x}{1+x}<\ln(1+x)<x.
\]

\hand{令 $f(t)=\ln(1+t)$，在区间 $[0,x]$ 上应用拉格朗日中值定理，得}
\[
\ln(1+x)=\frac{x}{1+\xi},\qquad \xi\in(0,x).
\]
\hand{由于 $0<\xi<x$，所以}
\[
\frac{1}{1+x}<\frac{1}{1+\xi}<1.
\]
\hand{两边同乘 $x>0$，得到}
\[
\frac{x}{1+x}<\ln(1+x)<x.
\]
\end{example}

\begin{example}
求极限
\[
\lim_{x\to\infty}x^2\left(\arctan\frac{1}{x}-\arctan\frac{1}{x+1}\right).
\]

\hand{令}
\[
f(t)=\arctan\frac{1}{t},
\]
\hand{则在区间 $[x,x+1]$ 上应用拉格朗日中值定理，有某个 $\xi\in(x,x+1)$，使}
\[
\arctan\frac{1}{x}-\arctan\frac{1}{x+1}=f(x)-f(x+1)=f'(\xi)(x-(x+1))=-f'(\xi).
\]
\hand{又因为}
\[
f'(t)=\frac{-1/t^2}{1+1/t^2}=-\frac{1}{1+t^2},
\]
\hand{所以原式等于}
\[
\lim_{x\to\infty}\frac{x^2}{1+\xi^2}=1.
\]
\end{example}

\paragraph{可以用拉格朗日中值定理证明的几个常识}
\subsubsection{推论 1}
若 $f(x)$ 在 $(a,b)$ 内可导，且
\[
f'(x)\equiv 0,
\]
则 $f(x)$ 在 $(a,b)$ 内恒为常数。

\hand{证明：任取 $x_1<x_2$，则 $f(x)$ 在 $[x_1,x_2]$ 上连续，在 $(x_1,x_2)$ 内可导。由拉格朗日中值定理，存在 $\xi\in(x_1,x_2)$，使}
\[
f(x_2)-f(x_1)=f'(\xi)(x_2-x_1)=0.
\]
\hand{故 $f(x_1)=f(x_2)$，从而 $f(x)$ 为常数。}

\subsubsection{推论 2}
若 $f(x),g(x)$ 在 $(a,b)$ 内可导，且
\[
f'(x)=g'(x),
\]
则在 $(a,b)$ 内
\[
f(x)=g(x)+C,
\]
其中 $C$ 为常数。

\hand{证明：令 $F(x)=f(x)-g(x)$，则}
\[
F'(x)=f'(x)-g'(x)=0.
\]
\hand{由推论 1 得 $F(x)=C$，即}
\[
f(x)=g(x)+C.
\]

\subsubsection{推论 3}
已知对任意 $x\in(a,b)$，有 $f'(x)>0\ (<0)$，则 $f(x)$ 在 $(a,b)$ 内单调递增（递减）。

\hand{证明：任取 $x_1<x_2$，由拉格朗日中值定理，存在 $\xi\in(x_1,x_2)$，使}
\[
f(x_2)-f(x_1)=f'(\xi)(x_2-x_1).
\]
\hand{若 $f'(x)>0$，则上式右端大于零，从而 $f(x_2)>f(x_1)$，故单调递增；$f'(x)<0$ 时同理。}

\begin{example}[3.5]
证明恒等式：
\[
\arcsin x+\arccos x=\frac{\pi}{2}\qquad (-1\le x\le 1).
\]

\hand{设}
\[
F(x)=\arcsin x+\arccos x.
\]
\hand{则 $F(x)$ 在 $[-1,1]$ 上连续，在 $(-1,1)$ 上可导，且}
\[
F'(x)=\frac{1}{\sqrt{1-x^2}}-\frac{1}{\sqrt{1-x^2}}=0.
\]
\hand{故 $F(x)$ 为常数。又}
\[
F(0)=\arcsin 0+\arccos 0=\frac{\pi}{2},
\]
\hand{所以}
\[
\arcsin x+\arccos x=\frac{\pi}{2}.
\]
\end{example}

\subsection{柯西中值定理}
设函数 $f(x)$ 和 $g(x)$ 满足
\begin{enumerate}[label=\arabic*)]
  \item 在闭区间 $[a,b]$ 上均连续；
  \item 在开区间 $(a,b)$ 内均可导；
  \item 对任意 $x\in(a,b)$，$g'(x)\neq 0$。
\end{enumerate}
则在 $(a,b)$ 内至少存在一点 $\xi\in(a,b)$，使得
\[
\frac{f(b)-f(a)}{g(b)-g(a)}=\frac{f'(\xi)}{g'(\xi)}.
\]

\begin{remark}
当 $g(x)=x$ 时，柯西中值定理即为拉格朗日中值定理。
\end{remark}

\begin{figure}[htbp]
\centering
\begin{tikzpicture}[scale=0.95, >=Latex]
  \draw[->] (0,0) -- (5.8,0) node[right] {$x$};
  \draw[->] (0,0) -- (0,4.1) node[above] {$y$};
  \draw[thick] (0.8,1.0) -- (4.9,2.2);
  \draw[thick, smooth] plot coordinates {(0.8,1.35) (1.6,1.55) (2.7,2.05) (3.7,2.95) (4.9,3.25)};
  \draw[dashed] (0.8,0) -- (0.8,1.0);
  \draw[dashed] (4.9,0) -- (4.9,2.2);
\end{tikzpicture}
\caption{柯西中值定理示意图}
\end{figure}

\begin{example}[3.6]
设 $0<a<b$，函数 $f(x)$ 在 $[a,b]$ 上连续，在 $(a,b)$ 内可导。证明存在一点 $\xi\in(a,b)$，使
\[
f(b)-f(a)=\xi f'(\xi)\ln\frac{b}{a}.
\]

\hand{取 $g(x)=\ln x$，由柯西中值定理得}
\[
\frac{f(b)-f(a)}{\ln b-\ln a}=\frac{f'(\xi)}{1/\xi}=\xi f'(\xi).
\]
\hand{于是}
\[
f(b)-f(a)=\xi f'(\xi)\ln\frac{b}{a}.
\]
\end{example}

\subsection{泰勒中值定理}
\paragraph{推导}
\hand{设}
\[
f(x)=a_0+a_1x+a_2x^2+\cdots+a_nx^n+o(x^n),
\]
\hand{在 $x=0$ 处比较各阶导数，得到}
\[
a_0=f(0),\quad a_1=f'(0),\quad a_2=\frac{f''(0)}{2},\quad \cdots,\quad a_n=\frac{f^{(n)}(0)}{n!}.
\]
\hand{因此}
\[
f(x)=f(0)+f'(0)x+\frac{f''(0)}{2!}x^2+\cdots+\frac{f^{(n)}(0)}{n!}x^n+o(x^n).
\]

\subsubsection{带皮亚诺余项型泰勒公式}
若函数 $f(x)$ 在包含 $x_0$ 处具有 $n$ 阶导数，那么在 $x_0$ 的一个邻域内，对该邻域内的任一 $x$，有
\[
f(x)=f(x_0)+f'(x_0)(x-x_0)+\frac{f''(x_0)}{2!}(x-x_0)^2+\cdots+\frac{f^{(n)}(x_0)}{n!}(x-x_0)^n+o[(x-x_0)^n].
\]

\hand{这里的 $o[(x-x_0)^n]$ 叫皮亚诺型余项。}
\hand{这种形式适合直接写函数在某点附近的展开式，也较适合求极限和做等价无穷小替换。}

\subsubsection{带拉格朗日型余项的泰勒公式}
若函数 $f(x)$ 在包含 $x_0$ 的某个区间 $(a,b)$ 内具有 $(n+1)$ 阶导数，则对任意 $x\in(a,b)$，有
\[
f(x)=f(x_0)+f'(x_0)(x-x_0)+\frac{f''(x_0)}{2!}(x-x_0)^2+\cdots+\frac{f^{(n)}(x_0)}{n!}(x-x_0)^n+\frac{f^{(n+1)}(\xi)}{(n+1)!}(x-x_0)^{n+1},
\]
其中 $\xi$ 是 $x$ 与 $x_0$ 之间的某个值。

\hand{皮亚诺余项适合直接写展开式，拉格朗日余项更适合证明题与误差估计。}
\hand{若写成 $f(x)=f(0)+f'(0)x+\dfrac{f''(0)}{2!}x^2+o(x^2)$，则提高一阶后可继续写成 $f(x)=f(0)+f'(0)x+\dfrac{f''(0)}{2!}x^2+\dfrac{f'''(0)}{3!}x^3+o(x^3)$。}

\begin{example}[3.8]
设
\[
\lim_{x\to 0}\frac{f(x)}{x}=1,\qquad f''(x)>0,
\]
证明：$f(x)>x$。

\hand{由极限可得 $f(0)=0,\ f'(0)=1$，故在 $x=0$ 附近有}
\[
f(x)=f(0)+f'(0)x+\frac{f''(\xi)}{2}x^2=x+\frac{f''(\xi)}{2}x^2>x.
\]
\hand{其中 $\xi$ 介于 $0$ 与 $x$ 之间，因为 $f''(\xi)>0$，所以二次项为正，从而 $f(x)>x$。}
\end{example}

\section{导数的微分学应用}

\subsection{单调性}
设函数 $f(x)$ 在闭区间 $[a,b]$ 上连续，在开区间 $(a,b)$ 内可导，则：
\begin{enumerate}[label=\arabic*)]
  \item 若在 $(a,b)$ 内 $f'(x)>0$（允许有限个点取零），则 $f(x)$ 在 $[a,b]$ 上单调增加；
  \item 若在 $(a,b)$ 内 $f'(x)<0$（允许有限个点取零），则 $f(x)$ 在 $[a,b]$ 上单调减少。
\end{enumerate}

\paragraph{求单调区间的一般步骤}
\begin{enumerate}[label=\arabic*)]
  \item 求导 $f'(x)$，找出驻点与不可导点；
  \item 以这些点为分界，讨论 $f'(x)$ 的符号；
  \item 根据导数符号写出函数的单调区间。
\end{enumerate}

\begin{example}
求函数
\[
y=e^{-x^2}
\]
的单调区间。
\end{example}

\begin{example}
设常数 $k>0$，函数
\[
f(x)=\ln x-\frac{k}{x}
\]
在 $(0,+\infty)$ 内零点个数为多少？
\end{example}

\subsection{极值}
\paragraph{求极值的基本步骤}
\begin{enumerate}[label=\arabic*)]
  \item 求 $f'(x)$，找驻点与不可导点；
  \item 用一阶导数符号法或二阶导数法判断极值；
  \item 作出结论。
\end{enumerate}

\subsection{最值}
\paragraph{闭区间上求最值的方法}
\begin{enumerate}[label=\arabic*)]
  \item 求导，找出区间内部的驻点与不可导点；
  \item 计算上述点与区间端点处的函数值；
  \item 比较这些值，得到最大值和最小值。
\end{enumerate}

\begin{example}
求函数
\[
y=x^2+2\cos x
\]
在区间 $[0,\pi]$ 上的最大值。
\end{example}

\begin{example}
将长为 $a$ 的铁丝截成两段，一段围成正方形，另一段围成圆，问怎样分配长度可以使总面积最小？
\end{example}

\subsection{曲线的凹凸性与拐点}
\paragraph{凹凸性的判定}
设函数 $f(x)$ 在 $[a,b]$ 上连续，在 $(a,b)$ 内具有二阶导数，则：
\begin{enumerate}[label=\arabic*)]
  \item 若在 $(a,b)$ 内 $f''(x)>0$，则曲线 $y=f(x)$ 在该区间上是凸的；
  \item 若在 $(a,b)$ 内 $f''(x)<0$，则曲线 $y=f(x)$ 在该区间上是凹的。
\end{enumerate}

\paragraph{拐点}
若曲线在某点两侧的凹凸性发生改变，则该点称为曲线的拐点。常见判别方法是考察 $f''(x)$ 的符号变化。

\begin{example}
问参数 $a,b$ 为何值时，点 $(1,3)$ 为曲线
\[
y=ax^3+bx^2
\]
的拐点。
\end{example}

\subsection{渐近线}
\paragraph{垂直渐近线}
若
\[
\lim_{x\to a^+}f(x)=\pm\infty
\quad \text{或} \quad
\lim_{x\to a^-}f(x)=\pm\infty,
\]
则直线 $x=a$ 是曲线 $y=f(x)$ 的一条垂直渐近线。

\paragraph{水平渐近线}
若
\[
\lim_{x\to +\infty}f(x)=b
\quad \text{或} \quad
\lim_{x\to -\infty}f(x)=b,
\]
则直线 $y=b$ 是曲线 $y=f(x)$ 的一条水平渐近线。

\paragraph{斜渐近线}
若存在常数 $a\neq 0,b$，使得
\[
\lim_{x\to \pm\infty}[f(x)-(ax+b)]=0,
\]
则直线 $y=ax+b$ 是曲线 $y=f(x)$ 的一条斜渐近线。

\begin{example}
讨论曲线
\[
y=\left(2-\frac{1}{x}\right)e^x
\]
的斜渐近线。
\end{example}

\section{整理说明}
这份 \LaTeX{} 文件按原 PDF 的顺序重排，标题编号也尽量对齐原讲义样式；原本的批注内容已经插回对应位置，并用不同字体区分。

\end{document}
